% -*- coding: utf-8 -*-

% \documentclass[12pt, a4paper, titlepage]{article}
\documentclass[language=en, year=1A, enterprise=ensIIE]{ensiie-report}

%%%%%%%%%%%%%%%%%%%%%%%%%%%%%%%%%%%%%%%%
%%%%%%% DÉCLARATION DU PRÉAMBULE %%%%%%%
%%%%%%%%%%%%%%%%%%%%%%%%%%%%%%%%%%%%%%%%

\usepackage{preamble}

%%%%%%%%%%%%%%%%%%%%%%%%%%%%%%%%%%%%%%%%
%%%%%%% DÉCLARATION DU GLOSSAIRE %%%%%%%
%%%%%%%%%%%%%%%%%%%%%%%%%%%%%%%%%%%%%%%%

\makeglossaries
% -*- coding: utf-8 -*-
% !TEX root = main.tex

%%%%%%%%%%%%%%%%%%%%%%%%%%%%%%
%%%%% DÉBUT DU GLOSSAIRE %%%%%
%%%%%%%%%%%%%%%%%%%%%%%%%%%%%%

%%%%%%%%%%%%%%%%%%%%%%
%%% MINIMUM REQUIS %%%
%%% - name         %%%
%%% - description  %%%
%%%%%%%%%%%%%%%%%%%%%%

\newglossaryentry{egal}
{
        name=égal,
        plural=égaux,
        feminine=égale,
        description={Qui a la même valeur, la même durée, la même importance, etc., que quelque chose d'autre}
}

\newglossaryentry{job}
{
        name=job,
        description={Unité de travail. Un job peut lancer plusieurs tâches}
}

%%%%%%%%%%%%%%%%%%%%%%%%%%%%%%
%%%%%% FIN DU GLOSSAIRE %%%%%%
%%%%%%%%%%%%%%%%%%%%%%%%%%%%%%

%%%%%%%%%%%%%%%%%%%%%%%%%%%%%%%
%%%%% DÉBUT DES ACRONYMES %%%%%
%%%%%%%%%%%%%%%%%%%%%%%%%%%%%%%

\newacronym{cpu}{CPU}{Central Processing Unit}

\newacronym{io}{I/O}{Input/Output}

\newacronym{go}{Go}{Gigaoctet}

%%%%%%%%%%%%%%%%%%%%%%%%%%%%%%%
%%%%%% FIN DES ACRONYMES %%%%%%
%%%%%%%%%%%%%%%%%%%%%%%%%%%%%%%

%%%%%%%%%%%%%%%%%%%%%%%%%%%%%%%
%%%%% DÉBUT DES SYMBOLES %%%%%%
%%%%%%%%%%%%%%%%%%%%%%%%%%%%%%%

\newglossaryentry{psi}
{%
  type=symbol,
  name={\ensuremath{\psi}},
  description={Fonction d'onde}
}

%%%%%%%%%%%%%%%%%%%%%%%%%%%%%%%
%%%%%% FIN DES SYMBOLES %%%%%%%
%%%%%%%%%%%%%%%%%%%%%%%%%%%%%%%


%%%%%%%%%%%%%%%%%%%%%%%%%%%%%%%%%%%%%%%%
% ACTIVATION DE L'EXTERNALISATION TIKZ %
%%%%%%%%%%%%%%%%%%%%%%%%%%%%%%%%%%%%%%%%

% On le fait pas dans le fichier .sty car c'est une commande qui dépend de documentclass.
% En clair, cela va permettre de créer les images TikZ dans le dossier "tikz/build" pour gagner du temps de compilation (si jamais on est amené à en faire)
% La première compilation est très longue, celles qui suivent sont théoriquement plus rapides.

% Externalisation intéressante SI vous créez un certain nombre de figures avec TikZ. Si ce n'est pas le cas, vous pouvez laisser commenter les lignes suivantes pour la désactiver.

% /!\ L'externalisation ne fonctionne pas sur Overleaf.

% \tikzexternalize
% [
%     prefix=tikz/build/,
%     up to date check={md5},
%     optimize command away=\AddToShipoutPictureBG
% ] % Chemin du dossier qui contient les graphiques TikZ

%%%%%%%%%%%%%%%%%%%%%%%%%%%%%%%%%%%%%
%%%%%%%%%%%% MÉTADONNÉES %%%%%%%%%%%%
%%%%%%%%%%%%%%%%%%%%%%%%%%%%%%%%%%%%%

% TODO
% Standard macros
\title{Analyse et Optimisation des performances du format de trace PALLAS pour les architectures de calcul exascale} % Sujet du stage
\author{Nikolaï MONTICELLI} % Prénom NOM de l'étudiant(e)
\date{17 juin 2025 to 29 août 2025} % Dates du stage

% TODO
% Custom macros
% À commenter si :
% - template utilisé pour autre chose qu'un rapport de stage
% - classe [ensiie-report] non utilisée
\years{2024-2025} % Année scolaire en cours
\logoSchool{img/logo_ensiie.png} % Chemin vers le logo de l'ensIIE (ex : img/logo_ensiie.png)
\logoEnterprise{img/LOGO_TSP.png} % Chemin vers le logo de l'entreprise du stage (ex : img/logo_entreprise.png)
\academicSupervisor{Dimitri WATEL} % Tuteur de stage
\internshipSupervisor{François TRAHAY} % Maître de stage
\internshipEnterprise{Télécom Sud Paris (Évry-Courcouronnes, 91)} % Lieu du stage

%%%%%%%%%%%%%%%%%%%%%%%%%%%%%%%%%%%%%%%
%%%%%%%%%% DÉBUT DU DOCUMENT %%%%%%%%%%
%%%%%%%%%%%%%%%%%%%%%%%%%%%%%%%%%%%%%%%

\begin{document}

\afterpage{\blankpage} % Page vide
\AddToShipoutPictureBG{\BackgroundPicFirstPage} % Pour activer le remplissage de la première page par le background choisi (dépend de l'entreprise définie dans la classe ensiie-report)
\maketitle
\ClearShipoutPictureBG % Pour retirer le fond sur toutes les pages suivantes

% \AddToShipoutPicture{\BackgroundPicAllPages} % Pour un autre fond sur toutes les autres pages

%%%%%%%%%%%%%%%%%%%%%%%%%%%%%%%%
%%%%%%%% REMERCIEMENTS %%%%%%%%%
%%%%%%%%%%%%%%%%%%%%%%%%%%%%%%%%

\pagenumbering{roman}
\pagestyle{FirstPages}

\section*{Remerciements}\label{sec:remerciements}
\addcontentsline{toc}{section}{Remerciements}

% -*- coding: utf-8 -*-
% !TEX root = ../main.tex

% TODO Vos remerciements (collègues de travail, famille, autre...)
H
\newpage

%%%%%%%%%%%%%%%%%%%%%%%%%%%%%%%%
%%%%%%%%%%% SYNTHÈSE %%%%%%%%%%%
%%%%%%%%%%%%%%%%%%%%%%%%%%%%%%%%

\section*{Synthèse}\label{sec:synthese}
\addcontentsline{toc}{section}{Synthèse}

% -*- coding: utf-8 -*-
% !TEX root = ../main.tex

% TODO Insérer la synthèse du stage

Pouvoir analyser les performances de son code de manière rapide et efficace est primordial pour les application parallèles développées à large échelle qui sont désormais essentielles dans de nombreux domaines comme la recherche, les simulations physiques, la prédiction de la météo, l'optimisation de la forme d'un aéronef, ... \newline
Le \verb!PEPR NumPEx! est un projet de recherche d'ampleur nationale dont l'objectif est la définition de la pile logicielle de la future machine de calcul exascale Alice Recoque qui sera installée au Très Grand Centre de Calcul du CEA de Bruyères-le-Châtel. Au sein de ce projet, l'équipe \verb!BENAGIL! du laboratoire \verb!Samovar! de \verb!Télécom SudParis! 
a pour objectif le développement d'une librairie, \verb!PALLAS!, qui permet l'analyse de la trace d'exécution d'un code qui soit adaptée au passage à l'échelle exascale. Ainsi, \verb!PALLAS! permet de générer une trace d'exécution contenant les détails de l'exéction d'une application. Ces traces, au format \verb!PALLAS! ont pour caractéristiques d'être de taille raisonnable, 
et d'avoir une incidence négligeable sur la durée d'exécution et les performances de l'application.\newline
Aujourd'hui, \verb!PALLAS! enregistre la trace et les métriques associées en mémoire pendant l'exécution de l'application en mémoire et ne les stocke sur le disque qu'à la fin de l'exécution. Cela peut-être problèmatique sur de nombreux aspects, notamment pour des traces assez volumineuses (plusieurs centaines de Gb).
L'objectif de ce stage de première année est alors d'analyser les performances de \verb!PALLAS!, et d'implémenter un flush périodique sur les disques suivant une métrique à définir pendant l'exécution même de l'applicaiton.


%%%
%%%

\paragraph*{Mots-clés}\label{par:mots-cles}

% TODO Mots-clés de votre rapport de stage

Trace d'exécution, performance, analyse, exascale

\newpage

%%%%%%%%%%%%%%%%%%%%%%%%%%%%%%%%
%%%%%%%%%% GLOSSAIRE %%%%%%%%%%%
%%%%%%%%%%%%%%%%%%%%%%%%%%%%%%%%

\section*{Glossaire}\label{sec:glossaire}
\addcontentsline{toc}{section}{Glossaire}

\setglossarystyle{list}
% \setupglossaries{numberedsection=autolabel}
% \setupglossaries{section,numberedsection=nameref}
\printglossary[type=\acronymtype, title=Acronymes, toctitle=Acronymes]
\printglossary[title=Vocabulaire, toctitle=Vocabulaire]
\printglossary[type=symbol, toctitle=Symboles]
\newpage

%%%%%%%%%%%%%%%%%%%%%%%%%%%%%%%%%%%%%%%%%%%
%%% LISTES DES FIGURES/TABLES/ÉQUATIONS %%%
%%%%%%%%%%%%%%%%%%%%%%%%%%%%%%%%%%%%%%%%%%%

\renewcommand{\listfigurename}{Liste des figures}
\listoffigures\label{sec:lof}
\addcontentsline{toc}{section}{Liste des figures}
% \renewcommand{\listtablename}{Liste des tables}
% \listoftables\label{sec:lot}
% \addcontentsline{toc}{section}{Liste des tables}
% \renewcommand{\lstlistlistingname}{Liste des listings}
% \lstlistoflistings\label{sec:lol}
% \addcontentsline{toc}{section}{Liste des listings}
\newpage

%%%%%%%%%%%%%%%%%%%%%%%%%%%%%%%%
%%%%%% TABLE DES MATIÈRES %%%%%%
%%%%%%%%%%%%%%%%%%%%%%%%%%%%%%%%

\renewcommand{\contentsname}{Table des matières}
\tableofcontents
\label{EndFirstPages} % Pour la numérotation en chiffre romain des premières pages
\newpage

%%%%%%%%%%%%%%%%%%%%%%%%%%%%%%%%
%%%%%%%%% INTRODUCTION %%%%%%%%%
%%%%%%%%%%%%%%%%%%%%%%%%%%%%%%%%

\pagenumbering{arabic}
\pagestyle{MainPages}

\section{Introduction}\label{sec:introduction}

% -*- coding: utf-8 -*-
% !TEX root = ../main.tex

\subsection{Télécom SudParis}\label{ssec:introduction_nom_entreprise}

% TODO Présenter l'entreprise et ses objectifs
Télécom SudParis est une grande école d'ingénieurs du numérique localisée à Évry, en France, et certifiée par la Commission des Titres d'Ingénieurs.\
Celle-ci fait partie de l'Institut Mines-Télécom et est membre de l'Institut Polytechnique de Paris.\\
L'école compte environ 1000 étudiants et plus d'une centaine de doctorants.\
Les thématiques étudiées tournent autour des Mathématiques Appliquées, de l'Informatique, des Réseaux, du Signal, et de la Physique.\\
L'équipe pédagogique est composée de plus d'une centaine d'enseignants chercheurs dont les travaux portent sur des thématiques comme les
réseaux complexes, le big data, l'IA, le cloud, l'Internet des objets ou encore la cybersécurité.\

% Exemple pour pour une image
% Rappel | Formats acceptés par \includegraphics (défaut) : png, jpg, pdf...
\begin{figure}[!h]
    \centering
    \includegraphics[width=0.7\textwidth]{img/campus_tsp.jpg}
    \caption{Photo du site}
    \label{fig:photo_site}
\end{figure}

%%%%%%%%%
%%%%%%%%%
%%%%%%%%%
%%%%%%%%%

\subsection{Service d'accueil}\label{ssec:introduction_service_accueil}

Ce stage a été encadré par \verb!François TRAHAY!, enseignant à \verb!Télécom SudParis! et chercheur au sein du laboratoire de recherche \verb!SAMOVAR! (Serivces répartis,\
Architectures, Modélisation, Validation, Administration des Réseux). Ce laboratoire accueille également quelques enseignants-chercheurs de l'\verb!ensIIE!, comme \verb!Valentin HONORÉ! qui a notamment co-encadré ce stage.\
Leurs travaux de recherche récents sont orientés vers l'analyse des performances dans le domaine du HPC (High Performance Computing).\
De plus, ils font partie de l'équipe \verb!BENAGIL! qui étudie la conception de systèmes distribués plus efficaces et plus sûrs en se focalisant sur leurs composants essentiels.

% TODO Présenter service/département de travail et son rôle dans entreprise
% On pourra aussi préciser le matériel à notre disposition pour mener notre stage

%%%%%%%%%
%%%%%%%%%
%%%%%%%%%
%%%%%%%%%

\subsection{Contexte et problématique}\label{ssec:introduction_contexte_problematique}

% TODO Présenter le contexte du stage et identifier les problèmes rencontrés qui justifient la pertinence du stage

%%%%%%%%%
%%%%%%%%%
%%%%%%%%%
%%%%%%%%%

\subsection{Objectifs du stage}\label{ssec:introduction_objectifs}

% TODO Présenter les objectifs du stage

%%%%%%%%%
%%%%%%%%%
%%%%%%%%%
%%%%%%%%%

\subsection{Contributions principales}\label{ssec:introduction_contributions}

% TODO Résumé des contributions majeures apportées

\newpage

%%%%%%%%%%%%%%%%%%%%%%%%%%%%%%%%
%%%%%%%%%%% CONTEXTE %%%%%%%%%%%
%%%%%%%%%%%%%%%%%%%%%%%%%%%%%%%%

\section{Contexte}\label{sec:contexte}

% -*- coding: utf-8 -*-
% !TEX root = ../main.tex

% TODO Présentation des travaux précédents pour expliquer d'où on part (programmes existants, études...)
% TODO Présenter les outils utilisés (logiciel/application, supercalculateur, langage de programmation...)
% Les sous-parties suivantes servent d'exemple, vous n'êtes pas obligés de les reprendre

\subsection{Programmes existants}\label{ssec:contexte_programmes}

\begin{figure}[!h]
\centering
\begin{minipage}{4cm}
\dirtree{%
.1 projet.
    .2 docs.
        .3 \( \ldots \).
    .2 src.
        .3 \( \ldots \).
    .2 tests.
        .3 \( \ldots \).
    .2 Makefile.
    .2 README.md.
    .2 \( \ldots \).
}
\end{minipage}
\caption{Exemple d'arborescence}
\label{fig:dirtree_example}
\end{figure}

%%%%%%%%%
%%%%%%%%%
%%%%%%%%%
%%%%%%%%%

\subsection{Outils utilisés}\label{ssec:outils}

\newpage

%%%%%%%%%%%%%%%%%%%%%%%%%%%%%%%%
%%%%%%% TRAVAUX RÉALISÉS %%%%%%%
%%%%%%%%%%%%%%%%%%%%%%%%%%%%%%%%

\section{Travaux réalisés}\label{sec:travaux_realises}

% -*- coding: utf-8 -*-
% !TEX root = ../main.tex

% TODO Développer dans plusieurs sous-sections ce que vous avez réalisé pendant votre stage

\tcbFileInput{shell}{\faTerminal}{lst/hw.sh}{Exemple de script shell}{lst:hw}

% Alternative | Fonctionne avec la version 2025 de TeX Live
% \fileInput[%
%     language=shell,
%     logo=\faTerminal,
%     caption={Exemple de script shell},
%     label=lst:hw
% ]{lst/hw.sh}

\tcbFileInput{python}{\faPython}{lst/dummy.py}{Exemple de script Python}{lst:dummy}

\begin{figure}[!h]
\begin{center}
\tikzsetnextfilename{dummy}
\resizebox{\textwidth}{!}{%
% -*- coding: utf-8 -*-
% !TEX root = ../main.tex

\documentclass[tikz]{standalone}

%%%%%%%%%%%%%%%%%%%%%%%%%%%%%%%%%%%%%%%%
%%%%%%% DÉCLARATION DU PRÉAMBULE %%%%%%%
%%%%%%%%%%%%%%%%%%%%%%%%%%%%%%%%%%%%%%%%

\makeatletter
\def\input@path{{../}}
\makeatother

\usepackage{preamble}

%%%%%%%%%%%%%%%%%%%%%%%%%%%%%%%%%%%%%%%
%%%%%%%%%% DÉBUT DU DOCUMENT %%%%%%%%%%
%%%%%%%%%%%%%%%%%%%%%%%%%%%%%%%%%%%%%%%

\begin{document}

%%%%%%%%%%%%%%%%%%%%%%%%%%%%%%%%%%%%%%%
%%%%%%% DÉBUT DE L'ILLUSTRATION %%%%%%%
%%%%%%%%%%%%%%%%%%%%%%%%%%%%%%%%%%%%%%%

\begin{tikzpicture}
\node(dummy) {Dummy};
\end{tikzpicture}

%%%%%%%%%%%%%%%%%%%%%%%%%%%%%%%%%%%%%%%
%%%%%%%% FIN DE L'ILLUSTRATION %%%%%%%%
%%%%%%%%%%%%%%%%%%%%%%%%%%%%%%%%%%%%%%%

\end{document}

}%
\caption{Exemple de dessin TikZ}
\label{fig:dummy}
\end{center}
\end{figure}

On peut insérer du texte avec une image \inlinegraphics{example-image-a} et faire à la \autoref{sec:introduction}.

Paralléliser un code Python, c'est une bonne idée. \cite{mpi4py}

Je fais référence au glossaire par l'adjectif suivant : \gls{egal}, \glsf{egal}, \glspl{egal}, \glsf{egal}s.

%%%%%%%%%
%%%%%%%%%
%%%%%%%%%
%%%%%%%%%

\subsection{Parlons \glsentrytext{cpu}}\label{ssec:dummy_cpu}

On parlera aussi éventuellement du \gls{cpu}, du débit \acrshort{io} ou de l'espace de stockage en \acrlong{go}.

Connaissez-vous la fonction \( \gls{psi} \) ?

\newpage

%%%%%%%%%%%%%%%%%%%%%%%%%%%%%%%%
%%%%%% RÉSULTATS OBTENUS %%%%%%%
%%%%%%%%%%%%%%%%%%%%%%%%%%%%%%%%

\section{Résultats obtenus}\label{sec:resultats_obtenus}

% -*- coding: utf-8 -*-
% !TEX root = ../main.tex

% TODO Présentation de vos résultats

\newpage

%%%%%%%%%%%%%%%%%%%%%%%%%%%%%%%%
%%%%%%%%%% CONCLUSION %%%%%%%%%%
%%%%%%%%%%%%%%%%%%%%%%%%%%%%%%%%

\section{Conclusion et perspectives}\label{sec:conclusion}

% -*- coding: utf-8 -*-
% !TEX root = ../main.tex

\subsection{Conclusion}\label{ssec:conclusion_conclusion}

%%%%%%%%%
%%%%%%%%%
%%%%%%%%%
%%%%%%%%%

\subsection{Retours d'expérience}\label{ssec:conclusion_retours}

% TODO Vos retours d'expérience : ce que vous avez appris/découvert

%%%%%%%%%
%%%%%%%%%
%%%%%%%%%
%%%%%%%%%

\subsection{Perspectives}\label{ssec:conclusion_perspectives}

% TODO Les perspectives d'amélioration ou d'utilisation de votre programme/de ce que vous avez développé

%%%%%%%%%
%%%%%%%%%
%%%%%%%%%
%%%%%%%%%

\newpage

%%%%%%%%%%%%%%%%%%%%%%%%%%%%%%%%
%%%%%%%%%%%%% DDRS %%%%%%%%%%%%%
%%%%%%%%%%%%%%%%%%%%%%%%%%%%%%%%

\appendix
\section{Développement durable et responsabilité sociétale}\label{sec:ddrs}

% -*- coding: utf-8 -*-
% !TEX root = ../main.tex

% TODO Annexe DDRS

%%%%%
%%%%%
%%%%%

\subsection{Développement durable}

Un exemple sympa. \cite{cea-terathermie}

%%%%%
%%%%%
%%%%%

\subsection{Responsabilité sociétale}

\newpage

%%%%%%%%%%%%%%%%%%%%%%%%%%%%%%%%%%%%
%%% ORGANISATION DE L'ENTREPRISE %%%
%%%%%%%%%%%%%%%%%%%%%%%%%%%%%%%%%%%%

\section{Organisation de l'entreprise}\label{sec:organisation}

% -*- coding: utf-8 -*-
% !TEX root = ../main.tex

% TODO Organigrammes de l'entreprise pour comprendre comment l'entreprise est structurée

\begin{center}
    \centering
    \includegraphics[height=0.915\textheight]{example-image-9x16}
    \label{fig:organigramme}
\end{center}

\newpage

%%%%%%%%%%%%%%%%%%%%%%%%%%%%%%%%
%%%%%% DÉTAILS TECHNIQUES %%%%%%
%%%%%%%%%%%%%%%%%%%%%%%%%%%%%%%%

\section{Détails techniques}\label{sec:details_techniques}

% -*- coding: utf-8 -*-
% !TEX root = ../main.tex

% TODO Expliquer la solution implémentée plus en détails (extraits de code, interfaces des applications, l’interface graphique des outils utilisés...)

\begin{terminalcode}
$ echo "Hello world!"
Hello world!
$ sleep 1
$
\end{terminalcode}

\newpage

%%%%%%%%%%%%%%%%%%%%%%%%%%%%%%%%
%%%%%%%%%% RÉFÉRENCES %%%%%%%%%%
%%%%%%%%%%%%%%%%%%%%%%%%%%%%%%%%

\renewcommand{\refname}{Références}
\bibliography{references}\label{sec:references}
\bibliographystyle{plain}
\newpage

%%%%%%%%%%%%%%%%%%%%%%%%%%%%%%%%
%%%%%%% FIN DU DOCUMENT %%%%%%%%
%%%%%%%%%%%%%%%%%%%%%%%%%%%%%%%%

\end{document}
